% !PRINTED 65
\weberpage{65}
Es ist aber nach (10)
\[
2a_2n_1n_2\omega_1 = a_1n_1n_2\omega_1 + a_2n_1n_2\omega_2 + mn_1n_2,
\]
und demnach wird diese letzte Formel:
\weq{16}{\begin{aligned}
&T(u+n_1\omega_1+n_2\omega_2)\\
&= e^{-\pi i[(n_1a_1+n_2a_2)(2u+n_1\omega_1+n_2\omega_2)+\pi i(b_1n_1+b_2n_2+mn_1n_2)]}T(u).
\end{aligned}}
% [sic] The printed exponent includes a nested \pi i factor in the final bracket.
Und hierin sind $n_1$ und $n_2$ beliebige ganze positive oder negative Zahlen.

Wenn man zwei $T$-Funktionen der Ordnung $m$ und $m'$ von den Charakteren $\gamma$, $\gamma'$ miteinander multipliziert, so entsteht eine neue $T$-Funktion, deren Ordnung $m+m'$ und deren Charakter $\gamma+\gamma'$ ist. Sind $(g_1, g_2)$ und $(g_1', g_2')$ die Charakteristiken der beiden Faktoren, so ist $(g_1+g_1', g_2+g_2')$ die Charakteristik des Produktes.

\begin{center}
\textbf{\S\,18. \gesperrt{Relationen zwischen verwandten $T$-Funktionen.}}
\end{center}

Zwei $T$-Funktionen von den gleichen Perioden, derselben Ordnung und demselben Charakter wollen wir \gesperrt{verwandt} nennen. Ist $T'(u)$ eine mit $T(u)$ verwandte $T$-Funktion und haben $a_1'$, $a_2'$, $b_1'$, $b_2'$ dieselbe Bedeutung für $T'$, wie $a_1$, $a_2$, $b_1$, $b_2$ für $T$, so ist infolge der vorausgesetzten Verwandtschaft [§\,17, (10), (11)]:
\weq{1}{a_1'\omega_2 - a_2'\omega_1 = a_1\omega_2 - a_2\omega_1,}
\weq{2}{b_1'\omega_2 - b_2'\omega_1 = b_1\omega_2 - b_2\omega_1 + 2n_1\omega_2 - 2n_2\omega_1,}
worin $n_1$, $n_2$ ganze Zahlen sind.

Daher lassen sich $\lambda$ und $\mu$ nach § 17, (9) so bestimmen, daß die Periodizitätsfaktoren der Funktion
\[
e^{-\pi i(\lambda u^2+\mu u)}T(u)
\]
dieselben werden wie die der Funktion $T'(u)$, und daher haben wir den Satz:

1. \gesperrt{Verwandten $T$-Funktionen können durch Hinzufügung einer $T$-Funktion nullter Ordnung:}
\[
L = Ce^{-\pi i(\lambda u^2+\mu u)}
\]
\gesperrt{als Faktor dieselben Periodizitätsfaktoren gegeben werden.}

Ferner folgt unmittelbar aus der Gleichheit der Charaktere:

2. \gesperrt{Wenn verwandte $T$-Funktionen $m$ter Ordnung $m-1$ gemeinsame Nullpunkte im Periodenparallelogramm haben, so haben sie auch den $m$ten gemeinsam;} und hieraus:
%CONT>>
